% Hình 3.1: Sơ đồ chuỗi thực hiện bộ ba [source / transition / follow-up] và kích hoạt Oracle 970
% Thiết kế đơn sắc (Black & White / Grayscale) cho in ấn học thuật, mở rộng khoảng cách dọc chống đè chữ.
\begin{tikzpicture}[
  >=Stealth,
  scale=0.88, transform shape,
  entity/.style={rectangle, draw=black, fill=gray!12, thick, minimum width=2.8cm, minimum height=0.9cm, align=center, rounded corners=2pt, font=\bfseries\small},
  server/.style={rectangle, draw=black, fill=gray!12, thick, minimum width=3.0cm, minimum height=0.9cm, align=center, rounded corners=2pt, font=\bfseries\small},
  db/.style={cylinder, draw=black, fill=gray!12, shape border rotate=90, aspect=0.25, thick, minimum width=2.4cm, minimum height=1.0cm, align=center, font=\bfseries\small},
  stepbadge/.style={font=\bfseries\footnotesize, align=right, anchor=east}
]

  % Các thực thể
  \node[entity] (fuzzer) at (0, 0)     {EvoMaster\\(bộ dựng và bộ phán quyết)};
  \node[server] (api)    at (5.4, 0)   {Máy chủ REST API\\(hệ thống mục tiêu)};
  \node[db]     (db)     at (10.8, 0)  {Cơ sở dữ liệu\\và kho phiên};

  % Đường gióng thời gian
  \draw[dashed, black!45] (fuzzer) -- (0, -12.2) coordinate (f_end);
  \draw[dashed, black!45] (api)    -- (5.4, -12.2) coordinate (a_end);
  \draw[dashed, black!45] (db)     -- (10.8, -12.2) coordinate (d_end);

  % ==================== 1. SOURCE ====================
  \node[stepbadge] at (-0.3, -1.2) {BƯỚC 1: SOURCE};
  \draw[->, thick] (0, -1.2) -- node[above, font=\footnotesize] {Lời gọi ghi, token $c_v$} (5.4, -1.2);
  \draw[->, thick] (5.4, -1.8) -- node[above, font=\scriptsize] {Ghi dữ liệu} (10.8, -1.8);
  \draw[<-, thick] (0, -2.4) -- node[above, font=\bfseries\footnotesize] {Phản hồi: \textbf{2xx}} (5.4, -2.4);

  % Lưu token (hộp chú thích độc lập)
  \node[rectangle, draw=black!70, fill=gray!6, rounded corners=2pt, font=\scriptsize, inner sep=3pt] at (2.7, -3.2) {
    Lưu lời gọi và token $c_v$ của nạn nhân
  };

  % ==================== 2. TRANSITION ====================
  \node[stepbadge] at (-0.3, -4.2) {BƯỚC 2: TRANSITION};
  \draw[->, thick] (0, -4.2) -- node[above, font=\footnotesize] {Hạ quyền nạn nhân, token $c_o$} (5.4, -4.2);
  \draw[->, thick] (5.4, -4.8) -- node[above, font=\scriptsize] {Cập nhật vai trò thấp hơn} (10.8, -4.8);
  \draw[<-, thick] (0, -5.4) -- node[above, font=\footnotesize] {Phản hồi: \textbf{2xx}} (5.4, -5.4);

  % ==================== 3. FOLLOW-UP ====================
  \node[stepbadge] at (-0.3, -6.6) {BƯỚC 3: FOLLOW-UP};
  \draw[->, thick] (0, -6.6) -- node[above, font=\footnotesize] {Gọi lại source, \textbf{token cũ} $c_v$} (5.4, -6.6);

  % ==================== 4. PHÁN QUYẾT ====================
  % Trường hợp A: 403 Forbidden (Hệ thống an toàn)
  \draw[<-, thick] (0, -7.6) -- node[above, font=\footnotesize] {Trường hợp A: \textbf{403}} (5.4, -7.6);
  \node[font=\footnotesize, align=center] at (2.7, -8.0) {
    $\implies$ phù hợp MR-A1
  };

  % Trường hợp B: 2xx OK (Nghi vấn lỗ hổng -> bước đăng nhập lại)
  \draw[<-, very thick] (0, -9.15) -- node[above, font=\bfseries\footnotesize] {Trường hợp B: \textbf{2xx} (dấu hiệu vi phạm)} (5.4, -9.15);
  
  \draw[->, dashed, thick] (0, -9.8) -- node[above, font=\scriptsize] {Đăng nhập lại, gửi bằng token mới $c_v'$} (5.4, -9.8);
  \draw[<-, dashed, thick] (0, -10.4) -- node[above, font=\scriptsize] {Token mới nhận \textbf{403}} (5.4, -10.4);

  \node[rectangle, draw=black, very thick, fill=gray!15, rounded corners=3pt, font=\bfseries\footnotesize, align=center, inner sep=4pt] at (2.7, -11.4) {
    Phát sinh cảnh báo Oracle 970
  };

\end{tikzpicture}
